\documentclass[a5paper,10pt,dvipdfmx]{jsarticle}
\usepackage{amsmath,amssymb}
\usepackage[margin=17mm,headsep=5mm,includehead]{geometry}
\usepackage{fancyhdr,needspace}
\usepackage[dvipdfmx,bookmarks=true,hidelinks]{hyperref}
\usepackage{pxjahyper}
\pagestyle{fancy}
\fancyhf{}
\fancyhead[L]{\small 数学原論}
\fancyhead[R]{\small 定義・命題一覧\quad\thepage}
\renewcommand{\headrulewidth}{0pt}
\setlength{\headheight}{14pt}
\setlength{\parindent}{0pt}
\setlength{\parskip}{0pt}
\setlength{\emergencystretch}{2em}
\raggedbottom
\newcommand{\entryhead}[2]{\par\addvspace{.55\baselineskip}\Needspace{4\baselineskip}\noindent\textbf{#1}\hfill{\small\ttfamily #2}\par\nobreak\smallskip}
\begin{document}
\begin{center}
{\small 数学原論}\par\vspace{3mm}
{\LARGE 定義・命題一覧}\par\vspace{2mm}
{\small 読者確認用}
\end{center}
知識台帳の検証済み項目（全248項目）を各ノートの掲載順に掲げる。証明・解説は収録しない。本文の厳密な記述については該当ノートを参照されたい。
\bigskip
\textbf{ノート別目次}\par\smallskip
\hyperlink{note-1}{N-1\quad 数学的言語}\dotfill 17項目\par
\hyperlink{note-2}{N-2\quad 集合・写像・関係}\dotfill 39項目\par
\hyperlink{note-3}{N-3\quad 選択公理}\dotfill 11項目\par
\hyperlink{note-4}{N-4\quad 順序と束}\dotfill 29項目\par
\hyperlink{note-5}{N-5\quad 自然数と有限集合}\dotfill 24項目\par
\hyperlink{note-6}{N-6\quad 整列と極大原理}\dotfill 12項目\par
\hyperlink{note-7}{N-7\quad 濃度と可算性}\dotfill 13項目\par
\hyperlink{note-8}{N-8\quad 位相と開集合}\dotfill 13項目\par
\hyperlink{note-9}{N-9\quad 二項演算とモノイド}\dotfill 19項目\par
\hyperlink{note-14}{N-14\quad フィルターと収束}\dotfill 8項目\par
\hyperlink{note-18}{N-18\quad 圏と函手}\dotfill 8項目\par
\hyperlink{note-24}{N-24\quad 集合環とσ-加法族}\dotfill 16項目\par
\hyperlink{note-53}{N-53\quad 近傍と閉包}\dotfill 14項目\par
\hyperlink{note-54}{N-54\quad 開基と部分基}\dotfill 9項目\par
\hyperlink{note-55}{N-55\quad 連続写像と同相}\dotfill 8項目\par
\hyperlink{note-62}{N-62\quad 群と部分群}\dotfill 8項目\par
\par\bigskip
\Needspace{12\baselineskip}
\hypertarget{note-1}{}
\section*{N-1\quad 数学的言語}
\pdfbookmark[1]{N-1 数学的言語}{bookmark-1}
\Needspace{9\baselineskip}
\subsection*{第1章\quad 第1節}
\entryhead{規約（言明と論理結合子）}{K-000001}
\(P,\allowbreak Q\) を意味の定まった数学的言明とし、\(\neg P\)、\(P\land Q\)、\(P\lor Q\)（包含的）、\(P\Rightarrow \allowbreak Q\)、\(P\Leftrightarrow \allowbreak Q\) を各自然言語の略記とする。同値は双方の含意を意味する。\par
\entryhead{規約（量化）}{K-000002}
仮定： \(P(x)\) は対象範囲の各対象について意味を持つ。\par
対象範囲を指定し、\(\forall x\) \(P(x)\) は任意の \(x\) について、\(\exists x\) \(P(x)\) は \(P\) を満たす \(x\) の存在、\(\exists !x\) \(P(x)\) はただ一つの存在を表す。\par
\entryhead{規約（対象の範囲と括弧）}{K-000003}
\(x=\allowbreak y\) は同一対象、\(x\in A\) は \(A\) の元であることを表す。括弧等で作用範囲を明示し、束縛変数を他変数と衝突しないよう改名する。\par
\Needspace{9\baselineskip}
\subsection*{第1章\quad 第2節}
\entryhead{規約（制限付き量化）}{K-000004}
仮定： 展開形では \(A\) の全元を含む量化範囲上で \(P(x)\) が意味を持つ。\par
\(\forall x\in A,\allowbreak P(x)\) は \(\forall x(x\in A\Rightarrow \allowbreak P(x))\)、\(\exists x\in A,\allowbreak P(x)\) は \(\exists x(x\in A\land P(x))\) の略記。\(P\) が \(A\) 上だけで意味を持つ場合は制限付き量化を自然言語として直接解釈する。\(A\) に元がないとき全称は成立、存在は不成立。\par
\entryhead{規約（その他の略記）}{K-000005}
\(x\notin A\) は \(\neg (x\in A)\)、\(x\ne y\) は \(\neg (x=\allowbreak y)\)。同種量化の連記 \(\forall x,\allowbreak y\in A\) は \(\forall x\in A,\allowbreak \forall y\in A\) の略記とし、存在量化も同様。\par
\entryhead{命題（一意存在）}{K-000006}
仮定： \(x,\allowbreak y\) の対象範囲は同一、\(y\) は衝突しない新しい変数。\par
\(\exists !x\) \(P(x)\) \(\Leftrightarrow \allowbreak \) \(\exists x(P(x)\land \forall y(P(y)\Rightarrow \allowbreak y=\allowbreak x))\)。\par
\Needspace{9\baselineskip}
\subsection*{第1章\quad 第3節}
\entryhead{背景原理（古典的推論）}{K-000007}
排中律 \(P\lor \neg P\)、矛盾律 \(\neg (P\land \neg P)\)、矛盾を導く仮定の否定、矛盾する仮定からの任意の結論を背景の推論原理として採用する。\par
\entryhead{命題（否定と含意）}{K-000008}
\(\neg \neg P\Leftrightarrow \allowbreak P\)；\((P\Rightarrow \allowbreak Q)\Leftrightarrow \allowbreak (\neg Q\Rightarrow \allowbreak \neg P)\)；\(\neg (P\Rightarrow \allowbreak Q)\Leftrightarrow \allowbreak (P\land \neg Q)\)。\par
\entryhead{命題（交換律・結合律・冪等律）}{K-000009}
連言・選言の交換律、結合律および冪等律 \(P\land P\Leftrightarrow \allowbreak P\)、\(P\lor P\Leftrightarrow \allowbreak P\)。\par
\entryhead{命題（分配律と吸収律）}{K-000010}
\(P\land (Q\lor R)\Leftrightarrow \allowbreak (P\land Q)\lor (P\land R)\)；\(P\lor (Q\land R)\Leftrightarrow \allowbreak (P\lor Q)\land (P\lor R)\)；\(P\land (P\lor Q)\Leftrightarrow \allowbreak P\)；\(P\lor (P\land Q)\Leftrightarrow \allowbreak P\)。\par
\entryhead{命題（De Morgan の法則）}{K-000011}
\(\neg (P\lor Q)\Leftrightarrow \allowbreak \neg P\land \neg Q\)；\(\neg (P\land Q)\Leftrightarrow \allowbreak \neg P\lor \neg Q\)。\par
\entryhead{命題（含意の合成）}{K-000012}
\(P\Rightarrow \allowbreak Q\) と \(Q\Rightarrow \allowbreak R\) から \(P\Rightarrow \allowbreak R\) が従う。同値についても推移性が成立する。\par
\entryhead{命題（量化の否定）}{K-000013}
仮定： 各式の量化範囲は同一。\par
\(\neg \exists x\) \(P(x)\Leftrightarrow \allowbreak \forall x\) \(\neg P(x)\)；\(\neg \forall x\) \(P(x)\Leftrightarrow \allowbreak \exists x\) \(\neg P(x)\)。\par
\entryhead{命題（量化の交換）}{K-000014}
仮定： \(x,\allowbreak y\) の対象範囲は互いに依存しない。\par
\(\forall x\forall y\) \(P(x,\allowbreak y)\Leftrightarrow \allowbreak \forall y\forall x\) \(P(x,\allowbreak y)\)；\(\exists x\exists y\) \(P(x,\allowbreak y)\Leftrightarrow \allowbreak \exists y\exists x\) \(P(x,\allowbreak y)\)。\par
\entryhead{命題（量化と論理結合子）}{K-000015}
仮定： \(P,\allowbreak Q\) の量化対象範囲は同一。\par
\(\forall x(P(x)\land Q(x))\Leftrightarrow \allowbreak (\forall x\) \(P(x))\land (\forall x\) \(Q(x))\)；\(\exists x(P(x)\lor Q(x))\Leftrightarrow \allowbreak (\exists x\) \(P(x))\lor (\exists x\) \(Q(x))\)。\par
\entryhead{命題（量化の順序に関する注意）}{K-000016}
仮定： \(x,\allowbreak y\) の対象範囲は互いに依存しない。\par
\(\exists y\forall x\) \(P(x,\allowbreak y)\Rightarrow \allowbreak \forall x\exists y\) \(P(x,\allowbreak y)\)。逆は一般に不成立。全称量化の選言への分配および存在量化の連言への分配は同値として一般に不成立（二対象の反例）。\par
\entryhead{規約（論証の方法）}{K-000017}
直接証明、対偶、背理法、場合分け、全称言明、存在言明、一意存在言明の論証における運用上の約束。\par
\par\bigskip
\Needspace{12\baselineskip}
\hypertarget{note-2}{}
\section*{N-2\quad 集合・写像・関係}
\pdfbookmark[1]{N-2 集合・写像・関係}{bookmark-2}
\Needspace{9\baselineskip}
\subsection*{第1章\quad 第1節}
\entryhead{公理（外延性）}{K-000018}
集合 \(A,\allowbreak B\) は任意の対象 \(x\) について \(x\in A\Leftrightarrow \allowbreak x\in B\) が成り立つとき等しい。逆は等号の意味から従う。\par
\entryhead{定義（部分集合）}{K-000019}
\(A\subseteq \allowbreak B\) \(\Leftrightarrow \allowbreak \) \(\forall x\in A,\allowbreak \) \(x\in B\)。\(A\subsetneq B\) \(\Leftrightarrow \allowbreak \) \(A\subseteq \allowbreak B\) かつ \(A\ne B\)。\par
\entryhead{命題（包含の性質）}{K-000020}
\(A\subseteq \allowbreak A\)；\(A\subseteq \allowbreak B\) かつ \(B\subseteq \allowbreak C\) なら \(A\subseteq \allowbreak C\)；\(A=\allowbreak B\) \(\Leftrightarrow \allowbreak \) \(A\subseteq \allowbreak B\) かつ \(B\subseteq \allowbreak A\)。\par
\entryhead{公理（空集合の存在）}{K-000021}
元を一つも持たない集合が存在する。外延性により一意であり \(\varnothing \) と記す。\par
\entryhead{命題（空集合の特徴）}{K-000022}
任意の集合 \(A\) について \(\varnothing \subseteq \allowbreak A\)。\(A\ne \varnothing \) \(\Leftrightarrow \allowbreak \) \(\exists x\in A\)。\par
\entryhead{公理（対集合の存在）}{K-000023}
任意の集合 \(a,\allowbreak b\) に対し \(x\in C\) \(\Leftrightarrow \allowbreak \) \((x=\allowbreak a\) または \(x=\allowbreak b)\) を満たす集合 \(C\) が存在し一意。\(\{a,\allowbreak b\}\) と記す。\(\{a\}=\allowbreak \{a,\allowbreak a\}\)。\par
\entryhead{公理（部分集合の指定）}{K-000024}
集合 \(A\) の各元について意味を持つ数学的条件 \(P(x)\) に対し \(x\in B\) \(\Leftrightarrow \allowbreak \) \(x\in A\) かつ \(P(x)\) を満たす集合 \(B\) が存在し一意。\(\{x\in A\) \(\mid \) \(P(x)\}\) と記す。\par
\entryhead{命題（全ての集合からなる集合は存在しない）}{K-000025}
すべての集合を元とする集合は存在しない。\par
\entryhead{公理（合併集合の存在）}{K-000026}
集合 \(C\) に対して \(x\in U\) \(\Leftrightarrow \allowbreak \) \(\exists A\in C,\allowbreak \) \(x\in A\) を満たす集合 \(U\) が存在し一意。\(\bigcup C\) と記す。\par
\entryhead{公理（冪集合の存在）}{K-000027}
集合 \(A\) に対し \(B\in P(A)\) \(\Leftrightarrow \allowbreak \) \(B\subseteq \allowbreak A\) を満たす集合 \(P(A)\) が存在し一意。\par
\entryhead{公理（一意に指定された元の集合）}{K-000028}
集合 \(A\) と各 \(x\in A\) に対してただ一つの集合 \(y\) を指定する条件 \(P(x,\allowbreak y)\) が与えられたとき、それらの \(y\) 全体からなる集合が存在する。\par
\Needspace{9\baselineskip}
\subsection*{第1章\quad 第2節}
\entryhead{定義（二集合の演算）}{K-000029}
\(A\cup B=\allowbreak \bigcup \{A,\allowbreak B\}\)、\(A\cap B=\allowbreak \{x\in A\) \(\mid \) \(x\in B\}\)、\(A\setminus B=\allowbreak \{x\in A\) \(\mid \) \(x\notin B\}\)。\(B\subseteq \allowbreak A\) のとき \(A\setminus B\) は \(A\) における \(B\) の補集合。\par
\entryhead{命題（集合演算の法則）}{K-000030}
合併・共通部分は交換・結合・冪等・分配・吸収律を満たす。\(A\cup \varnothing =\allowbreak A\)、\(A\cap \varnothing =\allowbreak \varnothing \)、\(A\setminus A=\allowbreak \varnothing \)。\(A\subseteq \allowbreak B\) \(\Leftrightarrow \allowbreak \) \(A\cap B=\allowbreak A\) \(\Leftrightarrow \allowbreak \) \(A\cup B=\allowbreak B\)。\par
\entryhead{命題（補集合の法則）}{K-000031}
\(A,\allowbreak B\subseteq \allowbreak E\) のとき \(A\cup (E\setminus A)=\allowbreak E\)、\(A\cap (E\setminus A)=\allowbreak \varnothing \)、\(E\setminus (E\setminus A)=\allowbreak A\)、\(E\setminus (A\cup B)=\allowbreak (E\setminus A)\cap (E\setminus B)\)、\(E\setminus (A\cap B)=\allowbreak (E\setminus A)\cup (E\setminus B)\)。\par
\entryhead{構成（順序対）}{K-000032}
\((a,\allowbreak b)=\allowbreak \{\{a\},\allowbreak \{a,\allowbreak b\}\}\) と構成する。\((a,\allowbreak b)=\allowbreak (c,\allowbreak d)\) \(\Leftrightarrow \allowbreak \) \(a=\allowbreak c\) かつ \(b=\allowbreak d\)。\par
\entryhead{構成（直積）}{K-000033}
\(A\times B=\allowbreak \{(a,\allowbreak b)\) \(\mid \) \(a\in A,\allowbreak b\in B\}\) は \(P(P(A\cup B))\) の部分集合として存在する。\(A\times B=\allowbreak \varnothing \) \(\Leftrightarrow \allowbreak \) \(A=\allowbreak \varnothing \) または \(B=\allowbreak \varnothing \)。\par
\Needspace{9\baselineskip}
\subsection*{第1章\quad 第3節}
\entryhead{定義（写像）}{K-000034}
写像 \(f:A\to B\) は始域 \(A\)、終域 \(B\)、グラフ \(G\subseteq \allowbreak A\times B\) を持ち、任意の \(a\in A\) にただ一つの \(b\in B\) が対応するもの。\(f(a)=\allowbreak b\) \(\Leftrightarrow \allowbreak \) \((a,\allowbreak b)\in G\)。相等は始域・終域・グラフの一致。\par
\entryhead{命題（写像の指定）}{K-000035}
\(A\) の各 \(a\) に一意に \(b\in B\) を指定する条件 \(P(a,\allowbreak b)\) は一意の写像 \(f:A\to B\) を定める。空集合から任意の \(B\) への写像は一意。\(A\ne \varnothing \) なら \(A\to \varnothing \) は存在しない。\par
\entryhead{構成（合成と恒等写像）}{K-000036}
\(f:A\to B\)、\(g:B\to C\) の合成 \(g\circ f:A\to C\) は \((g\circ f)(a)=\allowbreak g(f(a))\) で定まる。\(\mathrm{id}_{A}(a)=\allowbreak a\)。合成は結合的、恒等写像は単位。\par
\entryhead{定義（像と逆像）}{K-000037}
\(f:A\to B\)、\(S\subseteq \allowbreak A\)、\(T\subseteq \allowbreak B\) に対し \(f(S)=\allowbreak \{b\in B\) \(\mid \) \(\exists a\in S,\allowbreak f(a)=\allowbreak b\}\)、\(f^{-1}(T)=\allowbreak \{a\in A\) \(\mid \) \(f(a)\in T\}\)。\par
\entryhead{命題（像・逆像と集合演算）}{K-000038}
\(S,\allowbreak S'\subseteq \allowbreak A\)、\(T,\allowbreak T'\subseteq \allowbreak B\) に対し \(f(S\cup S')=\allowbreak f(S)\cup f(S')\)、\(f(S\cap S')\subseteq \allowbreak f(S)\cap f(S')\)、\(f^{-1}(T\cup T')=\allowbreak f^{-1}(T)\cup f^{-1}(T')\)、\(f^{-1}(T\cap T')=\allowbreak f^{-1}(T)\cap f^{-1}(T')\)、\(f^{-1}(B\setminus T)=\allowbreak A\setminus f^{-1}(T)\)。\(S\subseteq \allowbreak f^{-1}(f(S))\)、\(f(f^{-1}(T))=\allowbreak T\cap f(A)\)。\par
\entryhead{定義（単射・全射・全単射）}{K-000039}
\(f:A\to B\) が単射 \(\Leftrightarrow \allowbreak \) \(\forall a,\allowbreak a'\in A,\allowbreak \) \(f(a)=\allowbreak f(a')\Rightarrow \allowbreak a=\allowbreak a'\)。全射 \(\Leftrightarrow \allowbreak \) \(\forall b\in B,\allowbreak \exists a\in A,\allowbreak f(a)=\allowbreak b\) \(\Leftrightarrow \allowbreak \) \(f(A)=\allowbreak B\)。双方を全単射という。\par
\entryhead{命題（合成と単射・全射）}{K-000040}
\(f:A\to B\)、\(g:B\to C\) がともに単射なら \(g\circ f\) は単射、ともに全射なら全射。\(g\circ f\) が単射なら \(f\) は単射、全射なら \(g\) は全射。\par
\entryhead{命題（逆写像）}{K-000041}
\(f:A\to B\) が全単射 \(\Leftrightarrow \allowbreak \) \(g:B\to A\) で \(g\circ f=\allowbreak \mathrm{id}_{A}\)、\(f\circ g=\allowbreak \mathrm{id}_{B}\) を満たすものが存在する。その \(g\) は一意。\par
\entryhead{構成（制限・包含・終域の縮小）}{K-000042}
\(S\subseteq \allowbreak A\) なら包含 \(i:S\to A,\allowbreak \) \(i(s)=\allowbreak s\) と制限 \(f\mid _{S}:S\to B\) が存在。\(f(A)\subseteq \allowbreak T\subseteq \allowbreak B\) なら \(f\) を \(A\to T\) に終域縮小できる。\par
\entryhead{命題（直積の普遍性）}{K-000043}
射影 \(p:A\times B\to A\)、\(q:A\times B\to B\) に対し任意の \(u:X\to A,\allowbreak v:X\to B\) を満たす \(w:X\to A\times B\) で \(p\circ w=\allowbreak u,\allowbreak q\circ w=\allowbreak v\) が一意に存在する。\par
\Needspace{9\baselineskip}
\subsection*{第1章\quad 第4節}
\entryhead{定義（集合族）}{K-000044}
添字集合 \(I\) の各 \(i\) に集合 \(A_{i}\) が一意に与えられる対応を集合族 \((A_{i})_{i\in I}\) という。一意に指定された元の集合の原理により値の集合 \(C=\allowbreak \{A_{i}\) \(\mid \) \(i\in I\}\) が存在し、\(I\to C\) として表せる。\par
\entryhead{構成（集合族の合併と共通部分）}{K-000045}
集合族の合併 \(\bigcup _{i\in I}A_{i}=\allowbreak \bigcup \{A_{i}\mid i\in I\}\)。すべて \(A_{i}\subseteq \allowbreak E\) のとき \(E\) における共通部分は \(\{x\in E\mid \forall i\in I,\allowbreak x\in A_{i}\}\)。\(I\ne \varnothing \) なら共通部分は \(E\) に依存せず存在、\(I=\allowbreak \varnothing \) なら \(E\) に等しい。\par
\entryhead{命題（集合族と逆像・補集合）}{K-000046}
\(T_{i}\subseteq \allowbreak B\) と \(f:A\to B\) に対し \(f^{-1}(\bigcup T_{i})=\allowbreak \bigcup f^{-1}(T_{i})\)、\(f^{-1}(\bigcap _{B}\) \(T_{i})=\allowbreak \bigcap _{A}\) \(f^{-1}(T_{i})\)。\(B\setminus \bigcup T_{i}=\allowbreak \bigcap _{B}(B\setminus T_{i})\)、\(B\setminus \bigcap _{B}\) \(T_{i}=\allowbreak \bigcup (B\setminus T_{i})\)。空添字を含み、共通部分は指定母集合内。\par
\entryhead{構成（写像の集合と直積族）}{K-000047}
\(B^{A}\) は \(A\to B\) の写像全体。グラフを \(P(A\times B)\) 内で分離してから写像に対応させて存在を得る。集合族の積 \(\prod A_{i}=\allowbreak \{f\in U^{I}\) \(\mid \) \(\forall i,\allowbreak f(i)\in A_{i}\},\allowbreak \) \(U=\allowbreak \bigcup A_{i}\)。\(I=\allowbreak \varnothing \) なら積は一元集合、空の因子があれば空。\par
\entryhead{命題（集合族の直積の普遍性）}{K-000048}
\(P=\allowbreak \prod A_{i}\)、\(p_{i}:P\to A_{i},\allowbreak f\mapsto f(i)\)。集合 \(X\) と各 \(i\) に写像 \(g_{i}:X\to A_{i}\) が一意に指定される族に対し \(p_{i}\circ g=\allowbreak g_{i}\) を満たす \(g:X\to P\) が一意に存在。空添字を含む。\par
\entryhead{構成（非交和とその普遍性）}{K-000049}
\(D=\allowbreak \coprod A_{i}=\allowbreak \{(i,\allowbreak a)\in I\times U\mid a\in A_{i}\}\)、\(\iota _{i}:A_{i}\to D\)、\(a\mapsto (i,\allowbreak a)\)。任意の \(g_{i}:A_{i}\to X\) の族に対し \(g:D\to X\) で \(g\circ \iota _{i}=\allowbreak g_{i}\) が一意に存在。\par
\Needspace{9\baselineskip}
\subsection*{第1章\quad 第5節}
\entryhead{定義（関係）}{K-000050}
\(A\) から \(B\) への関係は \(A\times B\) の部分集合 \(R\)。\(a\) \(R\) \(b\) \(\Leftrightarrow \allowbreak \) \((a,\allowbreak b)\in R\)。\(A=\allowbreak B\) なら \(A\) 上の関係。\par
\entryhead{構成（逆関係と関係の合成）}{K-000051}
\(R\subseteq \allowbreak A\times B\) の逆関係 \(R^{-1}\subseteq \allowbreak B\times A\)、\(b\) \(R^{-1}\) \(a\) \(\Leftrightarrow \allowbreak \) \(a\) \(R\) \(b\)。\(S\subseteq \allowbreak B\times C\) の合成 \(S\circ R\subseteq \allowbreak A\times C\)、\(a(S\circ R)c\) \(\Leftrightarrow \allowbreak \) \(\exists b\in B,\allowbreak a\) \(R\) \(b\) \(\land \) \(b\) \(S\) \(c\)。合成は結合的。\par
\entryhead{定義（同値関係）}{K-000052}
\(A\) 上の関係 \(\sim \) が反射的 \(x\sim x\)、対称的 \(x\sim y\Rightarrow \allowbreak y\sim x\)、推移的 \(x\sim y\land y\sim z\Rightarrow \allowbreak x\sim z\) なら同値関係。\par
\entryhead{構成（同値類と商集合）}{K-000053}
\([x]=\allowbreak \{y\in A\mid x\sim y\}\)。\(x\sim y\) \(\Leftrightarrow \allowbreak \) \([x]=\allowbreak [y]\) \(\Leftrightarrow \allowbreak \) \([x]\cap [y]\ne \varnothing \)。\(A/\sim =\allowbreak \{[x]\mid x\in A\}\subseteq \allowbreak P(A)\)。標準全射 \(\pi :A\to A/\sim \)。\par
\entryhead{命題（同値関係と分割）}{K-000054}
分割 \(C\subseteq \allowbreak P(A)\) は各元が非空、相異なる元は交わらず、\(\bigcup C=\allowbreak A\)。\(A/\sim \) は分割。逆に \(x\sim y\) \(\Leftrightarrow \allowbreak \) \(\exists C\in \)分割\(,\allowbreak x,\allowbreak y\in C\) は同値関係で同値類の集合は元の分割。\par
\entryhead{命題（商集合の普遍性）}{K-000055}
\(\sim \) が \(A\) 上の同値関係、\(\pi :A\to A/\sim \)、\(f:A\to B\) とする。\(\exists !g:A/\sim \to B,\allowbreak \) \(g\circ \pi =\allowbreak f\) \(\Leftrightarrow \allowbreak \) \(\forall x,\allowbreak y\in A,\allowbreak x\sim y\Rightarrow \allowbreak f(x)=\allowbreak f(y)\)。\par
\entryhead{命題（写像の標準分解）}{K-000056}
\(f:A\to B\) による \(x\sim _{f}\) \(y\) \(\Leftrightarrow \allowbreak \) \(f(x)=\allowbreak f(y)\) は同値関係。標準全射 \(A\to A/\sim _{f}\)、全単射 \(A/\sim _{f}\to f(A)\)、包含 \(f(A)\to B\) の合成が \(f\)。\par
\par\bigskip
\Needspace{12\baselineskip}
\hypertarget{note-3}{}
\section*{N-3\quad 選択公理}
\pdfbookmark[1]{N-3 選択公理}{bookmark-3}
\Needspace{9\baselineskip}
\subsection*{第1章\quad 第1節}
\entryhead{定義（選択函数）}{K-000057}
集合族 \((A_{i})_{i\in I}\)、\(U=\allowbreak \bigcup A_{i}\) に対し、\(f:I\to U\) で全 \(i\) について \(f(i)\in A_{i}\) を満たすものを選択函数という。選択函数の集合は N-2 の直積 \(\prod A_{i}\) に等しい。\par
\entryhead{命題（一意選択）}{K-000058}
仮定： \(P(i,\allowbreak a)\) は各 \(i\) と \(a\in A_{i}\) について意味を持つ。\par
各 \(i\in I\) に対して \(A_{i}\) 内で条件 \(P(i,\allowbreak a)\) を満たす \(a\) が一意に存在するとき、全 \(i\) で \(P(i,\allowbreak f(i))\) を満たす選択函数 \(f\) が一意に存在する。選択公理は不要。\par
\entryhead{命題（二つの集合からの選択）}{K-000059}
\(I=\allowbreak \{i,\allowbreak j\}\)、各 \(A_{k}\) が非空なら選択函数が存在する。\(i=\allowbreak j\) の場合も含む。選択公理は不要。\par
\entryhead{公理（選択公理）}{K-000060}
任意の集合 \(I\) および各 \(A_{i}\ne \varnothing \) を満たす集合族 \((A_{i})_{i\in I}\) に対し、選択函数 \(f:I\to \bigcup A_{i}\) が存在する。\par
\entryhead{命題（直積の非空性）}{K-000061}
任意の集合族について \(\prod A_{i}\ne \varnothing \) \(\Leftrightarrow \allowbreak \) \(\forall i\in I,\allowbreak A_{i}\ne \varnothing \)。十分性は選択公理と同値な一般命題であり、必要性は選択公理不要。\par
\Needspace{9\baselineskip}
\subsection*{第1章\quad 第2節}
\entryhead{定義（右逆写像と左逆写像）}{K-000062}
\(f:A\to B\)、\(g:B\to A\) に対し \(f\circ g=\allowbreak \mathrm{id}_{B}\) なら \(g\) は右逆、\(g\circ f=\allowbreak \mathrm{id}_{A}\) なら左逆。右逆を持つ写像は全射、左逆を持つ写像は単射。\par
\entryhead{命題（全射の右逆）}{K-000063}
仮定： 右逆の存在の十分性には選択公理を使用。逆方向は仮定としての全射の右逆存在のみを使用。\par
選択公理の下で任意の全射 \(f:A\to B\) は右逆写像を持つ。逆に全ての全射が右逆を持つと仮定すれば選択公理が従う。\par
\entryhead{命題（単射の左逆）}{K-000064}
単射 \(f:A\to B\) に対し \(A\ne \varnothing \) なら左逆が存在する。\(A=\allowbreak \varnothing \) のとき左逆が存在する必要十分条件は \(B=\allowbreak \varnothing \)。選択公理は不要。\par
\entryhead{命題（関係の写像による実現）}{K-000065}
仮定： \(R\) の全域性。逆向きの証明では関係についての一般命題を仮定する。\par
任意の関係 \(R\subseteq \allowbreak A\times B\) が \(\forall a\in A,\allowbreak \exists b\in B,\allowbreak aRb\) を満たすなら、\(\forall a\in A,\allowbreak aRf(a)\) となる写像 \(f:A\to B\) が存在する。この一般命題は選択公理と同値。\par
\Needspace{9\baselineskip}
\subsection*{第1章\quad 第3節}
\entryhead{命題（商写像の切断と代表元）}{K-000066}
同値関係 \(\sim \) on \(A\) に対し \(\pi :A\to A/\sim \) の右逆 \(s\) が存在する。その像 \(T\) は各同値類とただ一つの元で交わる。逆にそのような \(T\) から \(s\) が一意に定まる。任意の分割の横断集合存在は選択公理と同値。\par
\entryhead{命題（全射の積）}{K-000067}
集合族 \((A_{i}),\allowbreak (B_{i})\)、全射 \(f_{i}:A_{i}\to B_{i}\) の族について \(F:\prod A_{i}\to \prod B_{i}\)、\(F(a)(i)=\allowbreak f_{i}(a(i))\) は全射。写像 \(F\) 自体の構成は選択公理不要、全射性に用いる。\par
\par\bigskip
\Needspace{12\baselineskip}
\hypertarget{note-4}{}
\section*{N-4\quad 順序と束}
\pdfbookmark[1]{N-4 順序と束}{bookmark-4}
\Needspace{9\baselineskip}
\subsection*{第1章\quad 第1節}
\entryhead{定義（前順序と半順序）}{K-000068}
集合 \(P\) 上の反射的・推移的関係を前順序といい、更に反対称的なら半順序という。任意の二元が比較可能な半順序を全順序という。順序集合は \(P\) と半順序の組。空集合も許す。\par
\entryhead{命題（狭義順序）}{K-000069}
半順序から \(x<y\) \(\Leftrightarrow \allowbreak \) \(x\leq \allowbreak y\land x\ne y\) と定めると非反射的かつ推移的。逆に非反射的推移的関係 \(<\) から \(x\leq \allowbreak y\) \(\Leftrightarrow \allowbreak \) \(x<y\lor x=\allowbreak y\) と定めると半順序となる。両構成は互いに逆。\par
\entryhead{命題（部分順序集合と双対）}{K-000070}
部分集合 \(S\subseteq \allowbreak P\) には制限した順序が入り、双対 \(P^{\mathrm{op}}\) は \(x\leq \allowbreak _{\mathrm{op}}\) \(y\) \(\Leftrightarrow \allowbreak \) \(y\leq \allowbreak x\) による順序集合。\(P,\allowbreak Q\) の直積には成分ごとの順序が入る。冪集合 \(P(E)\) は包含により順序集合。\par
\entryhead{命題（前順序の商）}{K-000071}
前順序 \(P\) に \(x\sim y\) \(\Leftrightarrow \allowbreak \) \(x\leq \allowbreak y\) かつ \(y\leq \allowbreak x\) と定めると同値関係。\(P/\sim \) 上の \([x]\leq \allowbreak [y]\) \(\Leftrightarrow \allowbreak \) \(x\leq \allowbreak y\) は代表元に依存せず半順序となる。\par
\Needspace{9\baselineskip}
\subsection*{第1章\quad 第2節}
\entryhead{定義（上界・下界と最大元・極大元）}{K-000072}
\(S\subseteq \allowbreak P\) の上界は \(\forall s\in S,\allowbreak s\leq \allowbreak u\) となる \(u\in P\)、下界は双対。\(S\) 内の上界は最大元、\(S\) 内の下界は最小元。\(m\in S\) が極大 \(\Leftrightarrow \allowbreak \) \(\forall s\in S,\allowbreak m\leq \allowbreak s\Rightarrow \allowbreak m=\allowbreak s\)、極小は双対。最大・最小元は存在すれば一意かつ極大・極小。\par
\entryhead{定義（上限と下限）}{K-000073}
\(S\subseteq \allowbreak P\) の上限 \(\sup _{P}\) \(S\) は \(S\) の上界全体の最小元、下限 \(\inf _{P}\) \(S\) は下界全体の最大元。存在すれば一意。空集合の上限は \(P\) の最小元、下限は最大元（存在する場合）。\par
\entryhead{命題（上限と下限の比較）}{K-000074}
\(S\subseteq \allowbreak T\subseteq \allowbreak P\) なら、両上限が存在する場合 \(\sup \) \(S\leq \allowbreak \sup \) \(T\)、両下限が存在する場合 \(\inf \) \(T\leq \allowbreak \inf \) \(S\)。\(U(T)\subseteq \allowbreak U(S)\)、\(L(T)\subseteq \allowbreak L(S)\)。\par
\entryhead{定義（鎖と有向集合）}{K-000075}
鎖は全順序となる部分集合（空を許す）。非空部分集合 \(D\) が任意の \(x,\allowbreak y\in D\) に共通上界 \(z\in D\) を持つとき上向き有向という。下向き有向は双対。非空鎖は上向き・下向き有向。全順序集合では極大元は最大元、極小元は最小元。\par
\Needspace{9\baselineskip}
\subsection*{第1章\quad 第3節}
\entryhead{定義（順序保存・順序埋込み・順序同型）}{K-000076}
\(f:P\to Q\) が順序保存 \(\Leftrightarrow \allowbreak \) \(x\leq \allowbreak y\Rightarrow \allowbreak f(x)\leq \allowbreak f(y)\)、順序埋込み \(\Leftrightarrow \allowbreak \) \(x\leq \allowbreak y\Leftrightarrow \allowbreak f(x)\leq \allowbreak f(y)\)、順序同型 \(\Leftrightarrow \allowbreak \) 全射な順序埋込み。順序埋込みは単射、順序同型の逆写像は順序保存。合成は各性質を保つ。\par
\entryhead{命題（全単射な順序保存写像）}{K-000077}
全単射な順序保存写像は一般に順序同型ではない。ただし始域が全順序なら順序同型である。\par
\entryhead{命題（順序同型と上限・下限）}{K-000078}
順序同型 \(f:P\to Q\) と \(S\subseteq \allowbreak P\) に対し、\(\sup _{P}\) \(S\) が存在する iff \(\sup _{Q}\) \(f(S)\) が存在し、後者は \(f(\sup _{P}\) \(S)\)。下限も同様。\par
\Needspace{9\baselineskip}
\subsection*{第1章\quad 第4節}
\entryhead{定義（束）}{K-000079}
空でない順序集合 \(L\) が任意の二元 \(x,\allowbreak y\) の上限・下限を持つとき束。\(x\lor y=\allowbreak \sup \{x,\allowbreak y\},\allowbreak \) \(x\land y=\allowbreak \inf \{x,\allowbreak y\}\)。全順序集合が非空なら束。冪集合は \(\lor =\allowbreak \cup \)、\(\land =\allowbreak \cap \) の束。\par
\entryhead{命題（束演算の諸法則）}{K-000080}
束の \(\lor ,\allowbreak \land \) は交換律・結合律・冪等律・吸収律を満たす。\(x\leq \allowbreak y\) \(\Leftrightarrow \allowbreak \) \(x\lor y=\allowbreak y\) \(\Leftrightarrow \allowbreak \) \(x\land y=\allowbreak x\)。\par
\entryhead{命題（演算による束の特徴づけ）}{K-000081}
非空集合 \(L\) に交換・結合・冪等・吸収律を満たす二つの二項演算 \(\lor ,\allowbreak \land \) があるなら、\(x\leq \allowbreak y\) \(\Leftrightarrow \allowbreak \) \(x\lor y=\allowbreak y\) と定めると半順序となり、元の演算は二元の上限・下限となる。\par
\entryhead{命題（束演算の単調性）}{K-000082}
\(x\leq \allowbreak x'\)、\(y\leq \allowbreak y'\) なら \(x\lor y\leq \allowbreak x'\lor y'\) および \(x\land y\leq \allowbreak x'\land y'\)。\par
\entryhead{定義（部分束と積束）}{K-000083}
非空部分集合 \(M\subseteq \allowbreak L\) が \(L\) の結び・交わりで閉じるとき部分束。制限順序の二項上限・下限は \(L\) の演算に一致。二つの束の直積は成分ごとの演算で束。制限順序で束であっても部分束とは限らない。\par
\entryhead{定義（束準同型と束同型）}{K-000084}
束間の写像が \(\lor \) と \(\land \) を保存するとき束準同型。束同型は逆写像を持ち双方が束準同型であるもの。束準同型は順序保存。全単射な束準同型は束同型、束同型と順序同型は同値。\par
\Needspace{9\baselineskip}
\subsection*{第1章\quad 第5節}
\entryhead{定義（有界束と完備束）}{K-000085}
最小元 0 と最大元 1 を持つ束を有界束という（\(0,\allowbreak 1\) は記号）。任意の部分集合 \(S\subseteq \allowbreak L\)（空含む）が上限・下限を持つ順序集合を完備束という。完備束は非空な有界束。\par
\entryhead{命題（完備性の片側判定）}{K-000086}
順序集合 \(L\) で任意の部分集合が上限を持つなら完備束。双対に任意の部分集合が下限を持つだけでもよい。\(\inf \) \(S=\allowbreak \sup \) \(L(S)\)、\(\sup \) \(S=\allowbreak \inf \) \(U(S)\)。\par
\entryhead{命題（冪集合と完備束の積）}{K-000087}
\(P(E)\) は完備束、\(\sup \) \(S=\allowbreak \bigcup S\)、\(\inf \) \(S=\allowbreak \{x\in E\mid \forall A\in S,\allowbreak x\in A\}\)。任意の完備束族 \((L_{i})\) の直積は成分ごとの順序で完備束。上限・下限も成分ごと。選択公理不要。\par
\Needspace{9\baselineskip}
\subsection*{第1章\quad 第6節}
\entryhead{定義（分配束）}{K-000088}
束 \(L\) で \(x\land (y\lor z)=\allowbreak (x\land y)\lor (x\land z)\) および \(x\lor (y\land z)=\allowbreak (x\lor y)\land (x\lor z)\) が成り立つとき分配束。二つの分配律のいずれかから他方が従う。 冪集合は分配束。5元の束 \(M=\allowbreak \{\varnothing ,\allowbreak \{a\},\allowbreak \{b\},\allowbreak \{c\},\allowbreak E\}\)（\(E=\allowbreak \{a,\allowbreak b,\allowbreak c\}\)、包含順序、\(a,\allowbreak b,\allowbreak c\) 相異なる）は分配束ではない。\par
\entryhead{定義（補元と Boolean 束）}{K-000089}
有界束 \(L\) の \(x\) の補元とは \(x\land y=\allowbreak 0_{L}\)、\(x\lor y=\allowbreak 1_{L}\) を満たす \(y\)。任意の元が補元を持つ有界分配束を Boolean 束と呼ぶ。冪集合は補集合を補元とする Boolean 束。退化一元束を許す。 上記 \(M\) は各元に補元が存在するが分配束ではなく補元は一意でない。\par
\entryhead{命題（補元の一意性と De Morgan の法則）}{K-000090}
Boolean 束で補元は一意（\(x^{c}\) と書く）。\((x^{c})^{c}=\allowbreak x\)、\(0^{c}=\allowbreak 1\)、\(1^{c}=\allowbreak 0\)、\((x\lor y)^{c}=\allowbreak x^{c}\land y^{c}\)、\((x\land y)^{c}=\allowbreak x^{c}\lor y^{c}\)。\(x\leq \allowbreak y\) iff \(y^{c}\leq \allowbreak x^{c}\)。\par
\Needspace{9\baselineskip}
\subsection*{第1章\quad 第7節}
\entryhead{定義（閉包作用素）}{K-000091}
順序集合 \(P\) 上の写像 \(c:P\to P\) が順序保存、\(x\leq \allowbreak c(x)\)、\(c(c(x))=\allowbreak c(x)\) を満たすとき閉包作用素。\(\operatorname{Fix}(c)=\allowbreak \{x\in P\mid c(x)=\allowbreak x\}\)。\(c(x)\) は \(x\) 以上の固定元の最小元。\par
\entryhead{命題（閉包作用素と完備束）}{K-000092}
完備束 \(L\) 上の閉包作用素 \(c\) の固定元集合 \(C\) は制限順序で完備束。\(S\subseteq \allowbreak C\) について \(\inf _{C}\) \(S=\allowbreak \inf _{L}\) \(S\)、\(\sup _{C}\) \(S=\allowbreak c(\sup _{L}\) \(S)\)。空の場合も成立。\par
\entryhead{定義（Galois 接続）}{K-000093}
順序保存写像 \(f:P\to Q\)、\(g:Q\to P\) が \(f(x)\leq \allowbreak y\) \(\Leftrightarrow \allowbreak \) \(x\leq \allowbreak g(y)\) を満たすとき Galois 接続（順序保存型）。\(f\) を左随伴、\(g\) を右随伴という。\(\mathrm{id}_{P}\leq \allowbreak g\circ f\)、\(f\circ g\leq \allowbreak \mathrm{id}_{Q}\)。互いに一意。\(g\circ f\) は閉包作用素。\par
\entryhead{命題（Galois 接続と上限・下限）}{K-000094}
Galois 接続 \((f,\allowbreak g)\) において \(S\subseteq \allowbreak P\) の上限が存在すれば \(f(\sup _{P}\) \(S)=\allowbreak \sup _{Q}\) \(f(S)\)（右辺も存在）。\(T\subseteq \allowbreak Q\) の下限について \(g(\inf _{Q}\) \(T)=\allowbreak \inf _{P}\) \(g(T)\)。\par
\entryhead{命題（像と逆像の Galois 接続）}{K-000095}
写像 \(h:A\to B\) に対する像写像 \(h_{\ast }:P(A)\to P(B)\) と逆像写像 \(h^{\ast }:P(B)\to P(A)\) は順序保存で \(h(S)\subseteq \allowbreak T\) iff \(S\subseteq \allowbreak h^{-1}(T)\)。合併と共通部分の保存が随伴から従う。\par
\entryhead{定理（完備束の不動点定理）}{K-000096}
完備束 \(L\) 上の順序保存写像 \(f:L\to L\) の固定元集合は制限順序で完備束。特に固定元が存在する。\(S\subseteq \allowbreak \operatorname{Fix}(f)\) の上限は \(\inf _{L}\{u\in L\) \(\mid \) \(\forall s\in S,\allowbreak s\leq \allowbreak u\) かつ \(f(u)\leq \allowbreak u\}\)、下限は双対。選択公理不要。\par
\par\bigskip
\Needspace{12\baselineskip}
\hypertarget{note-5}{}
\section*{N-5\quad 自然数と有限集合}
\pdfbookmark[1]{N-5 自然数と有限集合}{bookmark-5}
\Needspace{9\baselineskip}
\subsection*{第1章\quad 第1節}
\entryhead{定義（後者と帰納的集合）}{K-000097}
集合 \(x\) の後者 \(S(x)=\allowbreak x\cup \{x\}\) とする。\(0=\allowbreak \varnothing \) を含み，\(x\in E\Rightarrow \allowbreak S(x)\in E\) を満たす集合 \(E\) を帰納的集合という。\par
\entryhead{公理（無限公理）}{K-000098}
帰納的集合が存在する。\par
\entryhead{構成（自然数の集合）}{K-000099}
全ての帰納的集合に含まれる帰納的集合 \(\mathbb{N}\) が一意に存在する。これを自然数の集合とする。\par
\entryhead{定理（数学的帰納法）}{K-000100}
\(B\subseteq \allowbreak \mathbb{N}\) が \(0\in B\) および \(n\in B\Rightarrow \allowbreak S(n)\in B\) を満たすなら \(B=\allowbreak \mathbb{N}\)。\par
\entryhead{命題（自然数の基本性質）}{K-000101}
各 \(n\in \mathbb{N}\) は推移的な集合（\(x\in n\Rightarrow \allowbreak x\subseteq \allowbreak n\)）で，\(n\subseteq \allowbreak \mathbb{N}\)，\(n\notin n\)。各 \(x\in n\) は \(x\notin x\)。\par
\entryhead{命題（後者の単射性と前者）}{K-000102}
\(S:\mathbb{N}\to \mathbb{N}\) は単射で 0 はその像に属さない。0 以外の自然数は一意の自然数の後者である。\par
\entryhead{命題（自然数の順序）}{K-000103}
\(m<n\) \(\Leftrightarrow \allowbreak \) \(m\in n\)，\(m\leq \allowbreak n\) \(\Leftrightarrow \allowbreak \) \(m\subseteq \allowbreak n\) と定めると \(\mathbb{N}\) は全順序集合となる。\(m<S(n)\) \(\Leftrightarrow \allowbreak \) \(m\leq \allowbreak n\)。\par
\entryhead{命題（最小元の原理と強い帰納法）}{K-000104}
\(\mathbb{N}\) の非空部分集合は最小元を持つ。各 \(n\) について全ての \(m<n\) における \(P(m)\) から \(P(n)\) が従えば，全ての \(n\) で \(P(n)\)。\par
\Needspace{9\baselineskip}
\subsection*{第1章\quad 第2節}
\entryhead{定理（再帰定理）}{K-000105}
集合 \(X\)，\(a\in X\)，写像 \(F:\mathbb{N}\times X\to X\) に対し \(u:\mathbb{N}\to X\) で \(u(0)=\allowbreak a\)，\(u(S(n))=\allowbreak F(n,\allowbreak u(n))\) を満たすものが一意に存在する。\par
\entryhead{命題（帰納的体系の一意性）}{K-000106}
集合 \(X\) に \(a\in X\) と単射 \(t:X\to X\) があり，\(a\notin t(X)\)，\(a\) を含み \(t\) で閉じる部分集合が \(X\) だけなら，一意の全単射 \(u:\mathbb{N}\to X\) が \(u(0)=\allowbreak a\)，\(u(S(n))=\allowbreak t(u(n))\) を満たす。\par
\entryhead{定義（加法と乗法）}{K-000107}
自然数の加法・乗法を \(m+0=\allowbreak m\)，\(m+S(n)=\allowbreak S(m+n)\)，\(m\cdot 0=\allowbreak 0\)，\(m\cdot S(n)=\allowbreak m\cdot n+m\) で定める。\(1=\allowbreak S(0)\)。\par
\entryhead{命題（加法の法則）}{K-000108}
自然数の加法は結合的・可換で 0 を単位元とし，\(m+n=\allowbreak m+k\Rightarrow \allowbreak n=\allowbreak k\)。\(m\leq \allowbreak n\) \(\Leftrightarrow \allowbreak \) \(\exists k\in \mathbb{N},\allowbreak n=\allowbreak m+k\)。\par
\entryhead{命題（乗法の法則）}{K-000109}
自然数の乗法は結合的・可換で 1 を単位元とし，0 は吸収元。加法に対する分配律を満たす。\par
\Needspace{9\baselineskip}
\subsection*{第1章\quad 第3節}
\entryhead{定義（有限集合と濃度）}{K-000110}
集合 \(A\) がある \(n\in \mathbb{N}\) と全単射を持つとき有限集合という。\par
\entryhead{命題（有限集合の単射比較）}{K-000111}
自然数 \(m,\allowbreak n\) の間に単射 \(m\to n\) が存在する iff \(m\leq \allowbreak n\)。特に自然数との全単射が存在する場合，その自然数は一意。これを有限集合の元の個数 \(\mid A\mid \) とする。\par
\entryhead{命題（有限集合への一元の添加）}{K-000112}
有限集合 \(A\) と \(a\notin A\) に対し \(A\cup \{a\}\) は有限で \(\mid A\cup \{a\}\mid =\allowbreak S(\mid A\mid )\)。\par
\entryhead{命題（有限集合の部分集合と像）}{K-000113}
有限集合 \(A\) の部分集合 \(B\) は有限で \(\mid B\mid \leq \allowbreak \mid A\mid \)，等号は \(B=\allowbreak A\) のときに限る。任意の写像 \(f:A\to X\) の像は有限。\par
\entryhead{命題（有限集合間の写像）}{K-000114}
有限集合 \(A,\allowbreak B\) の間の単射の存在 iff \(\mid A\mid \leq \allowbreak \mid B\mid \)。全射 \(f:A\to B\) なら \(\mid B\mid \leq \allowbreak \mid A\mid \)。\(\mid A\mid =\allowbreak \mid B\mid \) のとき単射と全射と全単射は同値。\par
\entryhead{命題（有限集合の非交和）}{K-000115}
互いに交わらない有限集合 \(A,\allowbreak B\) の合併は有限で \(\mid A\cup B\mid =\allowbreak \mid A\mid +\mid B\mid \)。\par
\entryhead{命題（有限集合の直積・合併）}{K-000116}
有限集合 \(A,\allowbreak B\) の直積は有限で \(\mid A\times B\mid =\allowbreak \mid A\mid \cdot \mid B\mid \)。任意の二つの有限集合の合併も有限。\par
\entryhead{命題（有限集合の冪集合）}{K-000117}
有限集合 \(A\) の冪集合は有限である。\par
\Needspace{9\baselineskip}
\subsection*{第1章\quad 第4節}
\entryhead{定義（有限列と有限族）}{K-000118}
添字集合が自然数 \(n\) である集合族又は写像を長さ \(n\) の有限列と呼ぶ。有限集合を添字集合とする族を有限族と呼ぶ。\par
\entryhead{定理（有限選択）}{K-000119}
有限集合 \(I\) と各 \(A_{i}\ne \varnothing \) を満たす集合族に対し \(\prod A_{i}\ne \varnothing \)。選択公理を用いない。\par
\entryhead{命題（有限個の有限集合の合併と積）}{K-000120}
有限集合を添字集合とする有限集合の族の合併及び直積は有限。空の添字集合も含む。\par
\par\bigskip
\Needspace{12\baselineskip}
\hypertarget{note-6}{}
\section*{N-6\quad 整列と極大原理}
\pdfbookmark[1]{N-6 整列と極大原理}{bookmark-6}
\Needspace{9\baselineskip}
\subsection*{第1章\quad 第1節}
\entryhead{定義（整列集合と始切片）}{K-000121}
全順序集合 \(W\) の空でない任意の部分集合が最小元を持つとき整列集合という。\(W_{<a}=\allowbreak \{x\in W\) \(\mid \) \(x<a\}\) を \(a\) の真の始切片とする。始切片は下方閉な部分集合。\par
\entryhead{命題（始切片の形）}{K-000122}
整列集合 \(W\) の真の始切片は一意の \(a\in W\) によって \(W_{<a}\) と表される。任意の部分集合は制限順序で整列集合。\par
\entryhead{定理（超限帰納法）}{K-000123}
整列集合 \(W\) の各 \(a\) について，全 \(x<a\) における \(P(x)\) から \(P(a)\) が従うなら，全 \(a\in W\) について \(P(a)\)。\par
\entryhead{定理（超限再帰）}{K-000124}
整列集合 \(W\) と集合 \(X\) に対し，各 \(a\in W\) および \(h:W_{<a}\to X\) に値 \(F(a,\allowbreak h)\in X\) が一意に指定されるなら，\(u:W\to X\) で \(u(a)=\allowbreak F(a,\allowbreak u\mid _{W_{<a}})\) を満たすものが一意に存在する。\par
\entryhead{命題（整列集合の自己埋込み）}{K-000125}
整列集合 \(W\) の順序埋込み \(f:W\to W\) は全ての \(a\) で \(a\leq \allowbreak f(a)\)。従って \(W\) は自身の真の始切片と順序同型でない。\par
\entryhead{命題（整列集合の同型の一意性）}{K-000126}
整列集合 \(W,\allowbreak V\) の順序同型は存在すれば一意である。\(W\) の相異なる真の始切片は順序同型でない。\par
\entryhead{定理（整列集合の比較定理）}{K-000127}
任意の整列集合 \(W,\allowbreak V\) について，\(W\)と\(V\)が順序同型，\(W\)が\(V\)の真の始切片と順序同型，\(V\)が\(W\)の真の始切片と順序同型，のうちただ一つが成立する。\par
\Needspace{9\baselineskip}
\subsection*{第1章\quad 第2節}
\entryhead{定理（Hartogs の定理）}{K-000128}
任意の集合 \(A\) に対し，\(A\) への単射が存在しない整列集合 \(H\) が存在する。選択公理を用いない。\par
\Needspace{9\baselineskip}
\subsection*{第1章\quad 第3節}
\entryhead{定理（整列定理）}{K-000129}
仮定： 順方向は選択公理を使用。逆方向は整列可能性を仮定。\par
選択公理の下で任意の集合は整列順序を持つ。逆に任意の集合の整列可能性から選択公理が従う。\par
\Needspace{9\baselineskip}
\subsection*{第1章\quad 第4節}
\entryhead{定理（Hausdorff の極大原理）}{K-000130}
整列定理の下で任意の順序集合 \(P\) の任意の鎖 \(C\) は包含極大な鎖に含まれる。\par
\entryhead{定理（Zorn の補題）}{K-000131}
順序集合 \(P\) の任意の鎖（空の鎖を含む）が \(P\) 内に上界を持つなら \(P\) は極大元を持つ。選択公理の下で成立する。\par
\entryhead{定理（Zorn の補題からの選択公理）}{K-000132}
仮定： 逆方向は \(Zorn\) の補題を仮定し、選択公理を使用せず論証する。\par
\(Zorn\) の補題を一般命題として仮定すれば選択公理が従う。従って選択公理・整列定理・\(Hausdorff\) の極大原理・\(Zorn\) の補題は同値。\par
\par\bigskip
\Needspace{12\baselineskip}
\hypertarget{note-7}{}
\section*{N-7\quad 濃度と可算性}
\pdfbookmark[1]{N-7 濃度と可算性}{bookmark-7}
\Needspace{9\baselineskip}
\subsection*{第1章\quad 第1節}
\entryhead{定義（濃度の比較）}{K-000133}
\(A\)と\(B\)が全単射を持つとき同じ濃度という。\(\mid A\mid \leq \allowbreak \mid B\mid \) は単射 \(A\to B\) の存在，\(\mid A\mid =\allowbreak \mid B\mid \) は全単射の存在を表す。一般の濃度を集合として定義するものではない。\par
\entryhead{命題（濃度の比較の性質）}{K-000134}
同濃度は反射的・対称的・推移的。単射による比較は反射的・推移的。同濃度の集合で置換しても比較は不変。\par
\entryhead{定理（Cantor--Bernstein の定理）}{K-000135}
集合 \(A,\allowbreak B\) の間に単射 \(A\to B\) と \(B\to A\) が存在するなら全単射 \(A\to B\) が存在する。選択公理を用いない。\par
\entryhead{命題（全射と濃度の比較）}{K-000136}
単射 \(A\to B\) があり \(A\ne \varnothing \) なら全射 \(B\to A\) が存在する。\par
\entryhead{定理（Cantor の定理）}{K-000137}
任意の集合 \(A\) に対し \(\mid A\mid \leq \allowbreak \mid P(A)\mid \) だが，\(P(A)\) への全射 \(A\to P(A)\) は存在しない。従って \(\mid A\mid <\mid P(A)\mid \)。\par
\Needspace{9\baselineskip}
\subsection*{第1章\quad 第2節}
\entryhead{定義（高々可算・可算無限）}{K-000138}
集合 \(A\) が \(\mathbb{N}\) への単射を持つとき高々可算，\(\mathbb{N}\) と全単射を持つとき可算無限という。有限集合は高々可算。\par
\entryhead{命題（自然数集合の無限性）}{K-000139}
\(\mathbb{N}\) は有限集合でない。従って可算無限集合は有限でない。\par
\entryhead{定理（自然数の無限部分集合）}{K-000140}
\(\mathbb{N}\) の有限でない部分集合 \(B\) は \(\mathbb{N}\) と順序同型である。\par
\entryhead{命題（高々可算集合の分類）}{K-000141}
高々可算集合は有限集合又は可算無限集合のいずれか一方である。任意の部分集合及び像も高々可算。\par
\entryhead{定理（自然数の二項直積）}{K-000142}
\(\mathbb{N}\times \mathbb{N}\) は可算無限である。\par
\entryhead{命題（可算集合の有限積・有限合併）}{K-000143}
有限個の高々可算集合の直積及び合併は高々可算である。\par
\entryhead{命題（指定された単射を持つ可算族の合併）}{K-000144}
高々可算集合 \(I\) を添字とする集合族 \((A_{i})\) に，各単射 \(e_{i}:A_{i}\to \mathbb{N}\) が族として与えられていれば，合併は高々可算である。\par
\entryhead{系（冪集合の非可算性）}{K-000145}
\(P(\mathbb{N})\) は高々可算でない。\par
\par\bigskip
\Needspace{12\baselineskip}
\hypertarget{note-8}{}
\section*{N-8\quad 位相と開集合}
\pdfbookmark[1]{N-8 位相と開集合}{bookmark-8}
\Needspace{9\baselineskip}
\subsection*{第1章\quad 第1節}
\entryhead{定義（位相空間）}{K-000146}
集合 \(X\) の部分集合族 \(\mathcal{O}\subseteq \allowbreak P(X)\) が \(\varnothing ,\allowbreak X\in \mathcal{O}\)、有限交叉及び任意の合併で閉じるとき \(X\) 上の位相という。\((X,\allowbreak \mathcal{O})\) は位相空間、\(\mathcal{O}\) の元は開集合。空集合上の位相も許す。\par
\entryhead{命題（位相の有限交叉条件）}{K-000147}
\(\mathcal{O}\subseteq \allowbreak P(X)\) が位相であるためには \(X\in \mathcal{O}\)、二元の共通部分及び任意の合併で閉じることが必要十分。\par
\entryhead{定義（閉集合）}{K-000148}
位相空間 \((X,\allowbreak \mathcal{O})\) の部分集合 \(F\) は \(X\setminus F\in \mathcal{O}\) のとき閉集合である。\par
\entryhead{命題（閉集合系の特徴づけ）}{K-000149}
\(C\subseteq \allowbreak P(X)\) がある位相の閉集合全体となる iff \(\varnothing ,\allowbreak X\in C\)、任意交叉及び有限合併で閉じる。この位相は一意で \(\mathcal{O}=\allowbreak \{X\setminus F\) \(\mid \) \(F\in C\}\)。\par
\entryhead{定義（離散位相と密着位相）}{K-000150}
任意の集合 \(X\) 上で \(P(X)\) と \(\{\varnothing ,\allowbreak X\}\) は位相であり、それぞれ離散位相、密着位相という。 \(X\) が高々一元の場合は両者が一致する。\par
\entryhead{命題（二元集合上の位相）}{K-000151}
相異なる \(a,\allowbreak b\) に対して \(X=\allowbreak \{a,\allowbreak b\}\) 上の位相は \(\{\varnothing ,\allowbreak X\}\)、\(\{\varnothing ,\allowbreak \{a\},\allowbreak X\}\)、\(\{\varnothing ,\allowbreak \{b\},\allowbreak X\}\)、\(P(X)\) の四つである。\par
\Needspace{9\baselineskip}
\subsection*{第1章\quad 第2節}
\entryhead{定義（位相の粗密）}{K-000152}
同じ集合 \(X\) 上の位相 \(\mathcal{O},\allowbreak \mathcal{O}'\) に対し \(\mathcal{O}\subseteq \allowbreak \mathcal{O}'\) のとき \(\mathcal{O}\) は \(\mathcal{O}'\) より粗い、\(\mathcal{O}'\) は \(\mathcal{O}\) より細かいという。\par
\entryhead{命題（位相系の共通部分）}{K-000153}
\(X\) 上の位相の任意の族の共通部分は \(X\) 上の位相。空の族の場合は \(P(X)\)。\par
\entryhead{定理（生成位相）}{K-000154}
任意の \(S\subseteq \allowbreak P(X)\) に対し \(S\) を含む最小の位相 \(\mathcal{O}(S)\) が一意に存在し、\(S\) を含む全位相の共通部分で与えられる。\(S\subseteq \allowbreak T\Rightarrow \allowbreak \mathcal{O}(S)\subseteq \allowbreak \mathcal{O}(T)\)。\par
\entryhead{命題（位相全体の完備束）}{K-000155}
集合 \(X\) 上の位相全体は包含順序で完備束。任意の位相族の下限は共通部分、上限はその合併の生成位相。\par
\entryhead{命題（位相の合併は一般に位相でない）}{K-000156}
三元集合 \(X=\allowbreak \{a,\allowbreak b,\allowbreak c\}\) 上の二つの位相 \(\{\varnothing ,\allowbreak \{a\},\allowbreak X\}\) と \(\{\varnothing ,\allowbreak \{b\},\allowbreak X\}\) の合併は位相でない。\par
\Needspace{9\baselineskip}
\subsection*{第1章\quad 第3節}
\entryhead{命題（上集合の位相）}{K-000157}
順序集合 \((X,\allowbreak \leq \allowbreak )\) の上集合（\(x\in U\)、\(x\leq \allowbreak y\) なら \(y\in U\) となる部分集合 \(U\)）全体は位相で、任意交叉でも閉じる。\par
\entryhead{命題（有限集合上の位相）}{K-000158}
有限集合 \(X\) 上の位相は任意交叉でも閉じる。\par
\par\bigskip
\Needspace{12\baselineskip}
\hypertarget{note-9}{}
\section*{N-9\quad 二項演算とモノイド}
\pdfbookmark[1]{N-9 二項演算とモノイド}{bookmark-9}
\Needspace{9\baselineskip}
\subsection*{第1章\quad 第1節}
\entryhead{定義（二項演算）}{K-000159}
集合 \(M\) 上の二項演算とは写像 \(\mu :M\times M\to M\) である。通常 \(\mu (x,\allowbreak y)\) を \(xy\) と記す。\par
\entryhead{定義（結合律と交換律）}{K-000160}
二項演算がすべての \(x,\allowbreak y,\allowbreak z\in M\) に対し \((xy)z=\allowbreak x(yz)\) を満たすとき結合的、すべての \(x,\allowbreak y\in M\) に対し \(xy=\allowbreak yx\) を満たすとき可換という。\par
\entryhead{定義（単位元）}{K-000161}
左単位元 \(e\) はすべての \(x\in M\) について \(ex=\allowbreak x\)、右単位元 \(e\) は \(xe=\allowbreak x\) を満たす元。両方を満たす元を単位元という。\par
\entryhead{命題（単位元の一意性）}{K-000162}
左単位元 \(l\) と右単位元 \(r\) が存在すれば \(l=\allowbreak r\) で、その元は単位元。特に単位元は一意。\par
\entryhead{定義（半群とモノイド）}{K-000163}
非空集合 \(M\) と結合的二項演算の組を半群という。単位元を持つ半群をモノイドという。演算が可換なら可換半群、可換モノイドという。単位元は \(1_{M}\) と記す。 加法的記法では単位元を 0 と記す。\par
\entryhead{命題（モノイドの基本例）}{K-000164}
自然数集合 \(\mathbb{N}\) は加法及び乗法で可換モノイド。\(P(X)\) は合併及び共通部分で可換モノイド。\(X^{X}\) は写像の合成でモノイド。\par
\Needspace{9\baselineskip}
\subsection*{第1章\quad 第2節}
\entryhead{構成（有限列の積）}{K-000165}
モノイド \(M\) の有限列 \(x:n\to M\) の積 \(P(x)\) を、空列では \(1_{M}\)、\(n=\allowbreak S(m)\) のとき \(P(x)=\allowbreak P(x\mid _{m})\cdot x_{m}\) で一意に定める。\par
\entryhead{命題（有限列の連接と積）}{K-000166}
有限列 \(x,\allowbreak y\) の連接 \(x\ast y\) に対し \(P(x\ast y)=\allowbreak P(x)P(y)\)。空列も含む。\par
\entryhead{命題（有限積の括弧の独立性）}{K-000167}
空でない有限列の元を順序を変えず二項演算で任意に括弧づけして得た値は \(P(x)\) に等しい。\par
\entryhead{定義（自然数冪）}{K-000168}
\(x\in M\) の自然数冪 \(x^{n}\) は \(x\) の \(n\) 個の定値列の積。\(x^{0}=\allowbreak 1_{M}\)、\(x^{n+1}=\allowbreak x^{n}\) \(x\)。\par
\entryhead{命題（自然数冪の法則）}{K-000169}
\(m,\allowbreak n\in \mathbb{N}\) に対し \(x^{m+n}=\allowbreak x^{m}\) \(x^{n}\)、\((x^{m})^{n}=\allowbreak x^{mn}\)。\(xy=\allowbreak yx\) なら \((xy)^{n}=\allowbreak x^{n}y^{n}\)。\par
\Needspace{9\baselineskip}
\subsection*{第1章\quad 第3節}
\entryhead{定義（モノイド準同型）}{K-000170}
モノイド \(M,\allowbreak N\) の間の写像 \(f:M\to N\) が \(f(xy)=\allowbreak f(x)f(y)\)、\(f(1_{M})=\allowbreak 1_{N}\) を満たすとき準同型という。演算の保存のみでは単位元の保存は従わない。\par
\entryhead{命題（準同型と有限積）}{K-000171}
モノイド準同型 \(f\) は有限列の積を保存する。恒等写像と準同型の合成は準同型。\par
\entryhead{定義（モノイド同型）}{K-000172}
準同型 \(f:M\to N\) が逆写像を持ち、その逆写像も準同型であるときモノイド同型という。\par
\entryhead{命題（全単射なモノイド準同型）}{K-000173}
全単射なモノイド準同型はモノイド同型である。\par
\entryhead{定義（部分モノイド）}{K-000174}
モノイド \(M\) の部分集合 \(A\) が \(1_{M}\) を含み、積で閉じるとき部分モノイドという。制限された演算と同じ単位元でモノイドになる。\par
\entryhead{命題（準同型と部分モノイド）}{K-000175}
準同型 \(f:M\to N\) に対して部分モノイドの像及び逆像はそれぞれ部分モノイド。\par
\entryhead{定理（生成部分モノイド）}{K-000176}
部分集合 \(S\subseteq \allowbreak M\) を含む最小の部分モノイドは \(S\) の元の有限列の積全体（空積を含む）。 この集合を \(\langle S\rangle _{M}\) と記す。\par
\entryhead{命題（生成部分モノイドと準同型）}{K-000177}
準同型 \(f:M\to N\) と \(S\subseteq \allowbreak M\) に対して \(f(\langle S\rangle _{M})=\allowbreak \langle f(S)\rangle _{N}\)。\par
\par\bigskip
\Needspace{12\baselineskip}
\hypertarget{note-14}{}
\section*{N-14\quad フィルターと収束}
\pdfbookmark[1]{N-14 フィルターと収束}{bookmark-14}
\Needspace{9\baselineskip}
\subsection*{第1章\quad 第1節}
\entryhead{定義（フィルター）}{K-000233}
集合 \(X\) の部分集合族 \(F\) が \(X\in F\)、\(\varnothing \notin F\)、有限交叉と上方包含で閉じるときフィルターという。\par
\entryhead{命題（主フィルター）}{K-000234}
非空部分集合 \(A\subseteq \allowbreak X\) に対し \(\{B\subseteq \allowbreak X\mid A\subseteq \allowbreak B\}\) はフィルターである。\par
\Needspace{9\baselineskip}
\subsection*{第1章\quad 第2節}
\entryhead{定義（フィルター基）}{K-000235}
非空の部分集合族 \(B\subseteq \allowbreak P(X)\) が \(\varnothing \notin B\) かつ任意の \(A,\allowbreak B\in B\) に対し \(C\in B\)、\(C\subseteq \allowbreak A\cap B\) が存在するときフィルター基という。\par
\entryhead{命題（フィルター基の生成するフィルター）}{K-000236}
フィルター基 \(B\) に対し \(\{A\subseteq \allowbreak X\mid \)ある \(B\in B\) について \(B\subseteq \allowbreak A\}\) は \(B\) を含む最小のフィルター。\par
\Needspace{9\baselineskip}
\subsection*{第1章\quad 第3節}
\entryhead{命題（近傍フィルター）}{K-000237}
位相空間 \(X\) の各点 \(x\) の近傍系 \(\mathcal{N}(x)\) は \(X\) 上のフィルターである。\par
\entryhead{定義（フィルターの収束と集積点）}{K-000238}
位相空間 \(X\) のフィルター \(F\) が \(x\in X\) に収束するとは \(\mathcal{N}(x)\subseteq \allowbreak F\) をいう。\(x\) が \(F\) の集積点であるとは全ての \(A\in F\) に対し \(x\in \overline{A}\) となることをいう。\par
\entryhead{命題（収束と集積点の関係）}{K-000239}
フィルター \(F\) が \(x\) に収束するなら \(x\) は \(F\) の集積点である。\par
\entryhead{定理（集積点と細かいフィルター）}{K-000240}
フィルター \(F\) の集積点 \(x\) が存在する条件は、\(F\subseteq \allowbreak G\) かつ \(G\to x\) となるフィルター \(G\) が存在することと同値。\par
\par\bigskip
\Needspace{12\baselineskip}
\hypertarget{note-18}{}
\section*{N-18\quad 圏と函手}
\pdfbookmark[1]{N-18 圏と函手}{bookmark-18}
\Needspace{9\baselineskip}
\subsection*{第1章\quad 第1節}
\entryhead{定義（小圏と圏の規模）}{K-000241}
小圏 \(C\) は対象集合 \(Ob(C)\)、各対象 \(X,\allowbreak Y\) に対する射集合 \(Hom_{C}(X,\allowbreak Y)\)、合成写像と恒等射を持ち、結合律と単位律を満たす。大圏は対象全体を集合と仮定しない記述上の規約とする。\par
\entryhead{命題（集合とモノイドから得られる圏）}{K-000242}
集合と写像は大圏 \(Set\) をなし、任意のモノイド \(M\) は一対象の小圏を定める。\par
\Needspace{9\baselineskip}
\subsection*{第1章\quad 第2節}
\entryhead{定義（圏における同型射）}{K-000243}
射 \(f:X\to Y\) に対し \(g:Y\to X\) が \(g\circ f=\allowbreak 1_{X}\)、\(f\circ g=\allowbreak 1_{Y}\) を満たすとき \(f\) を同型射という。\par
\entryhead{命題（逆射の一意性）}{K-000244}
同型射の逆射は一意。恒等射は同型射であり同型射の合成は同型射。\par
\Needspace{9\baselineskip}
\subsection*{第1章\quad 第3節}
\entryhead{定義（函手）}{K-000245}
函手 \(F:C\to D\) は対象 \(X\) に \(F(X)\)、射 \(f:X\to Y\) に \(F(f):F(X)\to F(Y)\) を対応させ、恒等射と合成を保存する。\par
\entryhead{命題（函手の合成と恒等函手）}{K-000246}
恒等対応は函手であり、函手の合成は函手である。合成は結合律と単位律を満たす。\par
\entryhead{命題（函手による同型射の保存）}{K-000247}
函手 \(F\) は同型射を同型射に移し、逆射を逆射に移す。\par
\entryhead{命題（一対象圏とモノイド準同型）}{K-000248}
一対象圏の間の函手は対応するモノイド準同型と一対一に対応する。\par
\par\bigskip
\Needspace{12\baselineskip}
\hypertarget{note-24}{}
\section*{N-24\quad 集合環とσ-加法族}
\pdfbookmark[1]{N-24 集合環とσ-加法族}{bookmark-24}
\Needspace{9\baselineskip}
\subsection*{第1章\quad 第1節}
\entryhead{定義（集合環と集合代数）}{K-000178}
\(R\subseteq \allowbreak P(X)\) が \(\varnothing \) を含み、二元の合併及び差で閉じるとき集合環という。\(X\in R\) のとき集合代数という。\par
\entryhead{命題（集合環の有限演算）}{K-000179}
集合環は二元の共通部分及び有限合併で閉じる。空でない有限族の共通部分でも閉じる。\par
\entryhead{命題（集合代数の判定）}{K-000180}
\(A\subseteq \allowbreak P(X)\) が集合代数である iff \(X\in A\)、補集合及び二元の合併で閉じる。二元の共通部分で合併を置き換えても同値。\par
\entryhead{命題（集合環とBoolean束）}{K-000181}
集合環は包含順序で最小元 \(\varnothing \) を持つ分配束である。集合代数は最大元 \(X\) を持ち、補集合を補元とするBoolean束である。\par
\Needspace{9\baselineskip}
\subsection*{第1章\quad 第2節}
\entryhead{定義（σ-集合環とσ-加法族）}{K-000182}
集合環 \(R\) が任意の \(\mathbb{N}\) 添字列 \((A_{n})\) の合併 \(\bigcup _{n\in \mathbb{N}}\) \(A_{n}\) で閉じるとき \(\sigma -\)集合環という。\(X\in R\) なら \(\sigma -\)加法族（\(\sigma -\)代数）という。 組 \((X,\allowbreak A)\) は \(A\) が \(X\) 上の \(\sigma -\)加法族であるとき可測空間という。\par
\entryhead{命題（σ-集合環の可算演算）}{K-000183}
\(\sigma -\)集合環は高々可算の族の合併及び空でない高々可算族の共通部分で閉じる。\(\sigma -\)加法族では空の共通部分も許す。\par
\entryhead{命題（σ-加法族の判定）}{K-000184}
\(A\subseteq \allowbreak P(X)\) は \(\sigma -\)加法族 iff \(X\in A\)、補集合と \(\mathbb{N}\) 添字列の合併で閉じる。合併を共通部分で置き換えても同値。\par
\entryhead{命題（列の非交化）}{K-000185}
集合環 \(R\) の元の列 \((A_{n})\) に対し \(B_{n}=\allowbreak A_{n}\setminus \bigcup _{k<n}\) \(A_{k}\) とおくと \(B_{n}\in R\)、互いに素、\(\bigcup _{k\leq n}B_{k}=\allowbreak \bigcup _{k\leq n}A_{k}\)、\(\bigcup B_{n}=\allowbreak \bigcup A_{n}\)。\par
\entryhead{命題（互いに素な列による判定）}{K-000186}
集合環 \(R\) は互いに素な \(\mathbb{N}\) 添字列の合併で閉じる iff \(\sigma -\)集合環。\par
\Needspace{9\baselineskip}
\subsection*{第1章\quad 第3節}
\entryhead{定理（生成集合環・生成集合代数・生成σ-加法族）}{K-000187}
任意の \(S\subseteq \allowbreak P(X)\) を含む最小の集合環、集合代数、\(\sigma -\)集合環、\(\sigma -\)加法族がそれぞれ一意に存在する。各々 \(S\) を含む該当する全集合族の共通部分で与えられる。 順に \(r(S),\allowbreak a(S),\allowbreak r_{\sigma }(S),\allowbreak \sigma (S)\) と記す。\par
\entryhead{命題（生成作用素）}{K-000188}
\(r,\allowbreak a,\allowbreak r_{\sigma },\allowbreak \sigma \) は \(P(P(X))\) 上の閉包作用素である。生成される各種類の集合族全体は包含順序で完備束となる。\par
\entryhead{定理（生成集合代数の具体形）}{K-000189}
\(a(S)\) は \(S\) とその補集合からなる族の有限交叉の有限合併全体。空の交叉は \(X\)、空の合併は \(\varnothing \)。\par
\entryhead{命題（有限余有限集合代数）}{K-000190}
\(X=\allowbreak \mathbb{N}\times \{0,\allowbreak 1\}\) 上で有限集合又は補集合が有限の集合全体は集合代数だが \(\sigma -\)加法族でない。\par
\Needspace{9\baselineskip}
\subsection*{第1章\quad 第4節}
\entryhead{命題（逆像によるσ-加法族の判定）}{K-000191}
写像 \(f:X\to Y\)、\(X\) 上の \(\sigma -\)加法族 \(A\) に対し \(\{B\subseteq \allowbreak Y\) \(\mid \) \(f^{-1}(B)\in A\}\) は \(Y\) 上の \(\sigma -\)加法族。\par
\entryhead{命題（生成族による逆像条件）}{K-000192}
\(S\subseteq \allowbreak P(Y)\)、写像 \(f:X\to Y\)、\(X\) 上の \(\sigma -\)加法族 \(A\) について、すべての \(B\in S\) で \(f^{-1}(B)\in A\) ならすべての \(B\in \sigma (S)\) で \(f^{-1}(B)\in A\)。\par
\entryhead{命題（全射による逆像族）}{K-000193}
全射 \(f:X\to Y\) と \(Y\) 上の \(\sigma -\)加法族 \(B\) に対し \(\{f^{-1}(B)\mid B\in B\}\) は \(X\) 上の \(\sigma -\)加法族。\par
\par\bigskip
\Needspace{12\baselineskip}
\hypertarget{note-53}{}
\section*{N-53\quad 近傍と閉包}
\pdfbookmark[1]{N-53 近傍と閉包}{bookmark-53}
\Needspace{9\baselineskip}
\subsection*{第1章\quad 第1節}
\entryhead{定義（近傍と近傍系）}{K-000194}
位相空間 \((X,\allowbreak \mathcal{O})\) の点 \(x\in X\) に対し、\(x\in U\subseteq \allowbreak N\) となる \(U\in \mathcal{O}\) が存在する部分集合 \(N\subseteq \allowbreak X\) を \(x\) の近傍という。近傍全体を \(\mathcal{N}(x)\) と記す。\par
\entryhead{命題（近傍系の性質）}{K-000195}
位相空間 \(X\) の各点 \(x\) について、\((1)\) \(X\in \mathcal{N}(x)\)、\(N\in \mathcal{N}(x)\) なら \(x\in N\)、\((2)\) \(N\in \mathcal{N}(x)\) かつ \(N\subseteq \allowbreak M\subseteq \allowbreak X\) なら \(M\in \mathcal{N}(x)\)、\((3)\) \(N,\allowbreak M\in \mathcal{N}(x)\) なら \(N\cap M\in \mathcal{N}(x)\)、\((4)\) \(N\in \mathcal{N}(x)\) なら \(M\in \mathcal{N}(x)\) で全ての \(y\in M\) に対し \(N\in \mathcal{N}(y)\) となるものが存在する。\par
\entryhead{命題（開集合の近傍による判定）}{K-000196}
位相空間 \(X\) の部分集合 \(A\) が開集合であることと、全ての \(x\in A\) に対し \(A\in \mathcal{N}(x)\) となることは同値。\par
\entryhead{定理（近傍系による位相の特徴づけ）}{K-000197}
集合 \(X\) の各点 \(x\) に部分集合族 \(\mathcal{N}(x)\subseteq \allowbreak P(X)\) が与えられ、近傍系の四条件（K-000195）を満たすなら、それらを近傍系とする位相が一意に存在する。\par
\Needspace{9\baselineskip}
\subsection*{第1章\quad 第2節}
\entryhead{定義（内部と閉包）}{K-000198}
位相空間 \(X\) の部分集合 \(A\) の内部 \(\operatorname{Int}(A)\) は \(A\) に含まれる開集合全体の合併、閉包 \(\overline{A}\) は \(A\) を含む閉集合全体の共通部分。\par
\entryhead{命題（内部と閉包の極値性）}{K-000199}
位相空間 \(X\) の部分集合 \(A\) について、\(\operatorname{Int}(A)\) は \(A\) に含まれる最大の開集合、\(\overline{A}\) は \(A\) を含む最小の閉集合。\(A\) が開 \(\Leftrightarrow \allowbreak \) \(\operatorname{Int}(A)=\allowbreak A\)、\(A\) が閉 \(\Leftrightarrow \allowbreak \) \(\overline{A}=\allowbreak A\)。\par
\entryhead{命題（内部と閉包の補集合による双対性）}{K-000200}
位相空間 \(X\) の部分集合 \(A\) について、\(\operatorname{Int}(A)=\allowbreak X\setminus \overline{X\setminus A}\)、\(\overline{A}=\allowbreak X\setminus \operatorname{Int}(X\setminus A)\)。\par
\entryhead{命題（内部と閉包の近傍による特徴づけ）}{K-000201}
位相空間 \(X\) の \(A\subseteq \allowbreak X\)、\(x\in X\) に対して、\(x\in \operatorname{Int}(A)\) \(\Leftrightarrow \allowbreak \) \(A\in \mathcal{N}(x)\)。\(x\in \overline{A}\) \(\Leftrightarrow \allowbreak \) 全ての \(N\in \mathcal{N}(x)\) に対し \(N\cap A\ne \varnothing \)。\par
\entryhead{命題（内部と閉包の演算則）}{K-000202}
位相空間 \(X\) の部分集合 \(A,\allowbreak B\) について、内部と閉包は包含を保存し冪等である。\(\operatorname{Int}(A)\subseteq \allowbreak A\subseteq \allowbreak \overline{A}\)、\(\operatorname{Int}(X)=\allowbreak X\)、\(\overline{\varnothing }=\allowbreak \varnothing \)。\(\operatorname{Int}(A\cap B)=\allowbreak \operatorname{Int}(A)\cap \operatorname{Int}(B)\)、\(\overline{A\cup B}=\allowbreak \overline{A}\cup \overline{B}\)。\par
\Needspace{9\baselineskip}
\subsection*{第1章\quad 第3節}
\entryhead{定義（外部と境界）}{K-000203}
位相空間 \(X\) の部分集合 \(A\) の外部 \(\operatorname{Ext}(A)=\allowbreak \operatorname{Int}(X\setminus A)\)、境界 \(\partial A=\allowbreak \overline{A}\setminus \operatorname{Int}(A)\)。\par
\entryhead{命題（境界の特徴づけ）}{K-000204}
位相空間 \(X\) の部分集合 \(A\) について、\(\partial A=\allowbreak \overline{A}\cap \overline{X\setminus A}=\allowbreak \partial (X\setminus A)\) は閉集合。\(x\in \partial A\) \(\Leftrightarrow \allowbreak \) \(x\) の全ての近傍が \(A\) と \(X\setminus A\) の双方と交わる。\(X\) は \(\operatorname{Int}(A)\)、\(\operatorname{Ext}(A)\)、\(\partial A\) の互いに素な合併。\par
\entryhead{定義（集積点と導集合）}{K-000205}
位相空間 \(X\) の部分集合 \(A\) に対し、全ての \(N\in \mathcal{N}(x)\) が \(A\setminus \{x\}\) と交わる点 \(x\in X\) を \(A\) の集積点という。集積点全体を導集合 \(A'\) と記す。\par
\entryhead{命題（閉包と導集合）}{K-000206}
位相空間 \(X\) の部分集合 \(A\) について、\(\overline{A}=\allowbreak A\cup A'\)。\(A\) が閉集合 \(\Leftrightarrow \allowbreak \) \(A'\subseteq \allowbreak A\)。\par
\Needspace{9\baselineskip}
\subsection*{第1章\quad 第4節}
\entryhead{定理（Kuratowski の閉包公理）}{K-000207}
集合 \(X\) 上の写像 \(c:P(X)\to P(X)\) が \(c(\varnothing )=\allowbreak \varnothing \)、\(A\subseteq \allowbreak c(A)\)、\(c(c(A))=\allowbreak c(A)\)、\(c(A\cup B)=\allowbreak c(A)\cup c(B)\) を全ての \(A,\allowbreak B\subseteq \allowbreak X\) について満たすことと、\(c\) を閉包の写像とする位相が存在することは同値。その位相は一意。\par
\par\bigskip
\Needspace{12\baselineskip}
\hypertarget{note-54}{}
\section*{N-54\quad 開基と部分基}
\pdfbookmark[1]{N-54 開基と部分基}{bookmark-54}
\Needspace{9\baselineskip}
\subsection*{第1章\quad 第1節}
\entryhead{定義（開基）}{K-000208}
位相空間 \((X,\allowbreak \mathcal{O})\) の開集合族 \(\mathcal{B}\subseteq \allowbreak \mathcal{O}\) が開基であるとは、任意の開集合 \(U\) が \(\mathcal{B}\) の部分族の合併となることをいう。\par
\entryhead{命題（開基の局所判定）}{K-000209}
開集合族 \(\mathcal{B}\subseteq \allowbreak \mathcal{O}\) が開基であることと、任意の \(U\in \mathcal{O}\)、\(x\in U\) に対し \(x\in B\subseteq \allowbreak U\) となる \(B\in \mathcal{B}\) が存在することは同値。\par
\entryhead{定理（開基による位相の構成）}{K-000210}
集合 \(X\) の部分集合族 \(\mathcal{B}\) が或る位相の開基となる必要十分条件は、\(\bigcup \mathcal{B}=\allowbreak X\) であり、任意の \(B,\allowbreak C\in \mathcal{B}\) と \(x\in B\cap C\) に対し \(x\in D\subseteq \allowbreak B\cap C\) なる \(D\in \mathcal{B}\) が存在することである。その位相は一意。\par
\entryhead{命題（開基と位相の粗密）}{K-000211}
同じ集合 \(X\) 上の位相 \(\mathcal{O}_{1},\allowbreak \mathcal{O}_{2}\) の開基 \(\mathcal{B}_{1},\allowbreak \mathcal{B}_{2}\) に対し、\(\mathcal{O}_{1}\subseteq \allowbreak \mathcal{O}_{2}\) \(\Leftrightarrow \allowbreak \) 任意の \(B\in \mathcal{B}_{1}\) と \(x\in B\) に対し \(x\in C\subseteq \allowbreak B\) なる \(C\in \mathcal{B}_{2}\) が存在する。\par
\Needspace{9\baselineskip}
\subsection*{第1章\quad 第2節}
\entryhead{定義（部分基）}{K-000212}
位相空間 \((X,\allowbreak \mathcal{O})\) の部分集合族 \(\mathcal{S}\subseteq \allowbreak P(X)\) が部分基であるとは、\(\mathcal{S}\) の有限部分族の共通部分全体が開基であることをいう。空の共通部分は \(X\) とする。\par
\entryhead{定理（部分基による位相の生成）}{K-000213}
任意の集合 \(X\) と部分集合族 \(\mathcal{S}\subseteq \allowbreak P(X)\) に対し \(\mathcal{S}\) を部分基とする位相が一意に存在する。それは \(\mathcal{S}\) を含む最も粗い位相であり、\(\mathcal{S}\) の有限共通部分の任意合併全体である。\par
\Needspace{9\baselineskip}
\subsection*{第1章\quad 第3節}
\entryhead{定義（近傍基）}{K-000214}
位相空間 \(X\) の点 \(x\in X\) において、近傍族 \(\mathcal{B}_{x}\subseteq \allowbreak \mathcal{N}(x)\) が近傍基であるとは、任意の \(N\in \mathcal{N}(x)\) がある \(B\in \mathcal{B}_{x}\) を含むことをいう。\par
\entryhead{命題（開基から得られる近傍基）}{K-000215}
開基 \(\mathcal{B}\) に対し各点 \(x\) の族 \(\mathcal{B}_{x}=\allowbreak \{B\in \mathcal{B}\) \(\mid \) \(x\in B\}\) は \(x\) の近傍基である。\par
\entryhead{定理（近傍基からの位相の構成）}{K-000216}
集合 \(X\) の各点 \(x\) に非空の部分集合族 \(\mathcal{B}_{x}\) が与えられるとする。各 \(B\in \mathcal{B}_{x}\) は \(x\) を含む。任意の \(B,\allowbreak C\in \mathcal{B}_{x}\) に対して \(D\in \mathcal{B}_{x}\)、\(D\subseteq \allowbreak B\cap C\) が存在する。各 \(B\in \mathcal{B}_{x}\) に対して \(C\in \mathcal{B}_{x}\) が存在し、全ての \(y\in C\) に対して \(D\in \mathcal{B}_{y}\)、\(D\subseteq \allowbreak B\) が存在する。この三条件の下で、各 \(\mathcal{B}_{x}\) を近傍基とする位相が一意に存在する。逆も成立。\par
\par\bigskip
\Needspace{12\baselineskip}
\hypertarget{note-55}{}
\section*{N-55\quad 連続写像と同相}
\pdfbookmark[1]{N-55 連続写像と同相}{bookmark-55}
\Needspace{9\baselineskip}
\subsection*{第1章\quad 第1節}
\entryhead{定義（連続写像）}{K-000217}
位相空間 \(X,\allowbreak Y\) の写像 \(f:X\to Y\) が連続であるとは、\(Y\) の任意の開集合 \(V\) に対し \(f^{-1}(V)\) が \(X\) で開であることをいう。\par
\entryhead{命題（連続写像の合成）}{K-000218}
恒等写像は連続であり、連続写像の合成は連続である。\par
\entryhead{命題（閉集合による連続性判定）}{K-000219}
写像 \(f:X\to Y\) が連続 \(\Leftrightarrow \allowbreak \) 任意の閉集合 \(F\subseteq \allowbreak Y\) の逆像 \(f^{-1}(F)\) が閉集合。\par
\Needspace{9\baselineskip}
\subsection*{第1章\quad 第2節}
\entryhead{定義（点における連続性）}{K-000220}
写像 \(f:X\to Y\) が点 \(x\in X\) で連続であるとは、\(f(x)\) の任意の近傍 \(N\) に対し \(f^{-1}(N)\) が \(x\) の近傍であることをいう。\par
\entryhead{定理（連続性の局所判定）}{K-000221}
写像 \(f:X\to Y\) が連続 \(\Leftrightarrow \allowbreak \) \(X\) のすべての点で連続。\par
\entryhead{命題（閉包による連続性判定）}{K-000222}
写像 \(f:X\to Y\) が連続 \(\Leftrightarrow \allowbreak \) 任意の \(A\subseteq \allowbreak X\) について \(f(\overline{A})\subseteq \allowbreak \overline{f(A)}\)。\par
\Needspace{9\baselineskip}
\subsection*{第1章\quad 第3節}
\entryhead{定義（同相写像）}{K-000223}
連続な全単射 \(f:X\to Y\) で逆写像も連続なものを同相写像という。同相写像が存在するとき \(X,\allowbreak Y\) は同相という。\par
\entryhead{命題（同相写像と開集合）}{K-000224}
全単射 \(f:X\to Y\) が同相写像 \(\Leftrightarrow \allowbreak \) 任意の \(U\subseteq \allowbreak X\) について \(U\) が開であることと \(f(U)\) が開であることが同値。\par
\par\bigskip
\Needspace{12\baselineskip}
\hypertarget{note-62}{}
\section*{N-62\quad 群と部分群}
\pdfbookmark[1]{N-62 群と部分群}{bookmark-62}
\Needspace{9\baselineskip}
\subsection*{第1章\quad 第1節}
\entryhead{定義（群）}{K-000225}
モノイド \(G\) の各元 \(x\) に対して \(xy=\allowbreak yx=\allowbreak 1_{G}\) となる \(y\in G\) が存在するとき \(G\) を群という。積が交換的なら可換群という。\par
\entryhead{命題（逆元の一意性）}{K-000226}
群 \(G\) の各元 \(x\) の逆元は一意であり \(x^{-1}\) と記す。\((x^{-1})^{-1}=\allowbreak x\)、\((xy)^{-1}=\allowbreak y^{-1}x^{-1}\)。\par
\entryhead{命題（簡約律）}{K-000227}
群 \(G\) で \(ax=\allowbreak ay\) なら \(x=\allowbreak y\)、\(xa=\allowbreak ya\) なら \(x=\allowbreak y\)。方程式 \(ax=\allowbreak b\)、\(ya=\allowbreak b\) の解はそれぞれ \(a^{-1}b\)、\(ba^{-1}\) に限る。\par
\Needspace{9\baselineskip}
\subsection*{第1章\quad 第2節}
\entryhead{定義（部分群）}{K-000228}
群 \(G\) の部分集合 \(H\) が \(1_{G}\) を含み、積と逆元で閉じているとき部分群といい \(H\leq \allowbreak G\) と記す。\par
\entryhead{命題（部分群判定）}{K-000229}
部分集合 \(H\subseteq \allowbreak G\) が部分群 \(\Leftrightarrow \allowbreak \) \(H\ne \varnothing \) かつ任意の \(x,\allowbreak y\in H\) に対して \(xy^{-1}\in H\)。\par
\entryhead{命題（部分群の共通部分）}{K-000230}
群 \(G\) の部分群からなる任意の族の共通部分は部分群。空の族の場合は \(G\) とする。\par
\Needspace{9\baselineskip}
\subsection*{第1章\quad 第3節}
\entryhead{定義（生成部分群）}{K-000231}
部分集合 \(S\subseteq \allowbreak G\) を含む部分群全体の共通部分を生成部分群 \(\langle S\rangle \) という。これは \(S\) を含む最小の部分群。\par
\entryhead{命題（生成部分群の有限積表示）}{K-000232}
生成部分群 \(\langle S\rangle \) は \(S\cup \{s^{-1}\mid s\in S\}\) の元の有限積全体に等しい。空積は \(1_{G}\)。\par
\end{document}
